\chapter{Organização Curricular}

\textbf{FUNÇÃO 1:} Compreender os fundamentos científicos e tecnológicos da Inteligência Artificial para aplicação em soluções computacionais.

\textbf{FUNÇÃO 2:} Desenvolver sistemas computacionais baseados em Inteligência
Artificial, Análise Preditiva e Ecossistemas IoT, atendendo às normas de qualidade
corporativas, robustez metodológica, governança de dados e arquitetura segura.




\section{Concepção da Organização Curricular}
A organização curricular foi estruturada com base no desenvolvimento progressivo de competências profissionais, adotando uma abordagem por itinerário formativo.
A matriz curricular integra fundamentos científicos, competências tecnológicas, desenvolvimento de software, ciência de dados, inteligência artificial, computação distribuída e aplicações industriais, garantindo que cada unidade curricular contribua para a formação de um profissional capaz de atuar em projetos completos de Inteligência Artificial.
Ao longo do curso, o estudante será incentivado a desenvolver projetos integradores que consolidem as competências adquiridas, utilizando ferramentas, linguagens, frameworks e plataformas amplamente adotadas pelo mercado.

\section{Estrutura Curricular}

\begin{table}[htb]
\centering
\renewcommand{\arraystretch}{1.2} % Melhora o espaçamento entre as linhas
\begin{tabular}{|c|l|c|}
\hline
\textbf{Módulo} & \textbf{Unidade Curricular} & \textbf{Carga Horária} \\
\hline
Básico I & Fundamentos da Inteligência Artificial & 80 h \\
\hline
Básico I & Programação para Inteligência Artificial & 120 h \\
\hline
Específico I & Engenharia de Dados para IA & 100 h \\
\hline
Específico I & Machine Learning & 120 h \\
\hline
Específico I & Deep Learning e Redes Neurais & 120 h \\
\hline
Específico I & Inteligência Artificial Generativa e LLMs & 180 h \\
\hline
Específico II & Agentes Inteligentes, MCP, LangGraph e Sistemas Autônomos & 160 h \\
\hline
Específico II & Desenvolvimento de Software Inteligente com IA & 140 h \\
\hline
Específico II & Visão Computacional & 140 h \\
\hline
Específico III & IoT, IIoT e Computação Distribuída & 120 h \\
\hline
Específico III & MLOps, DevOps para IA e Observabilidade & 100 h \\
\hline
Específico III & Projeto Integrador & 160 h \\
\hline
\end{tabular}
% \caption{Estrutura Curricular}
% \label{tab:estrutura_curricular}
\end{table}

Carga Horária Total: \textbf{1.540 horas}

\section{Metodologia de Ensino}
O curso adota metodologias ativas de aprendizagem, privilegiando a experimentação prática e o desenvolvimento de competências profissionais.

Serão utilizadas as seguintes estratégias:
\begin{itemize}
    \item Aprendizagem Baseada em Projetos.
    \item Aprendizagem Baseada em Problemas.
    \item Estudos de Caso.
    \item Desafios Técnicos.
    \item Hackathons.
    \item Laboratórios Práticos.
    \item Desenvolvimento Colaborativo.
    \item Pair Programming.
    \item Code Review.
    \item Seminários Técnicos.
    \item Oficinas de Integração.
    \item Desenvolvimento de Projetos Reais.
    \item Mentorias Técnicas.
    \item Avaliações Formativas.
\end{itemize}

\section{Projeto Integrador}

\begin{itemize}
    \item Engenharia de Dados.
    \item Machine Learning ou Deep Learning.
    \item IA Generativa.
    \item Agentes Inteligentes.
    \item RAG.
    \item APIs.
    \item Banco Vetorial.
    \item Banco Relacional.
    \item Interface Web.
    \item Integração com serviços externos.
    \item Docker.
    \item Kubernetes.
    \item Pipeline CI/CD.
    \item MLOps.
    \item Observabilidade.
    \item Documentação Técnica.
    \item Apresentação Executiva.
\end{itemize}

A solução deverá ser apresentada a uma banca avaliadora composta por docentes e profissionais convidados, considerando critérios técnicos, inovação, aplicabilidade, documentação e qualidade da implementação.
